% from ega4-17-fr.tex
\oldpage[IV-4]{106}

les images des $(g_i)_x\in\sh{O}_{X,x}$ dans $\sh{O}_{X_s,x}=\sh{O}_{X,x}/\mathfrak m_s\sh{O}_{X,x}$ soient égales aux $t_i$. Soit $X'$ le sous-préschéma fermé de $V$ défini par l'idéal de $A$ engendré par les $g_i$; la suite $(t_i)$ étant par hypothèse régulière (0, 16.5.7) il résulte de (11.3.8) qu'en remplaçant au besoin $V$ par un voisinage ouvert plus petit, on peut supposer que le morphisme $X'\to S$, restriction de $f$, est plat et de présentation finie. D'autre part, puisque les $t_i$ forment un système de paramètres de $\sh{O}_{X_s,x}$, l'anneau $\sh{O}_{X'_s,x}$ est par définition artinien, et comme $x$ est fermé dans $X'_s$, on en conclut qu'il est isolé dans $X'_s$. En remplaçant $V$ au besoin par un voisinage plus petit de $x$, on en conclut, grâce à (13.1.4) que le morphisme $X'\to S$ est quasi-fini.

\begin{corollary}[17.16.2]
\label{IV.17.16.2}
Soit $f:X\to S$ un morphisme fidèlement plat et localement de présentation finie. Alors il existe un morphisme $g:S'\to S$, fidèlement plat, localement de présentation finie et localement quasi-fini, tel qu'il existe un $S$-morphisme $S'\to X$ (autrement dit, tel qu'il existe une $S'$-section de $X'=X\times_S S'$ (I, 3.3.14)). Si $S$ est quasi-compact (resp. quasi-compact et quasi-séparé), on peut supposer $S'$ affine (resp. $S'$ affine et le morphisme $g$ quasi-fini).
\end{corollary}

Pour tout $s\in S$, la fibre $X_s$ est non vide par hypothèse et est un préschéma localement de type fini sur $k(s)$; l'ensemble $U_s$ des points $x$ de $X_s$ où $\sh{O}_{X_s,x}$ est un anneau de Cohen--Macaulay est ouvert dans $X_s$ (6.11.3) et est non vide, puisqu'il contient les points maximaux de $X_s$ (0, 16.5.1); il contient par suite un point $x_s$ fermé dans $X_s$ (10.4.7). Soit $X'(s)$ un sous-préschéma affine de $X$ contenant $x_s$ et vérifiant les conditions de (17.16.1). Pour obtenir un préschéma $S'$ vérifiant les conditions de l'énoncé, il suffit de prendre la somme des $X'(s)$, où $s$ parcourt $S$. Puisque le morphisme $X'(s)\to S$ est plat et localement de présentation finie, l'image $U(s)$ de $X'(s)$ est ouverte dans $S$ (2.4.6); lorsque $S$ est quasi-compact, on peut donc recouvrir $S$ par un nombre fini de $U(s_i)$ et le préschéma $S'$ somme des $X'(s_i)$ répond encore à la question et est affine. Si de plus $S$ est quasi-séparé, on peut supposer les immersions ouvertes $U(s_i)\to S$ quasi-compactes (I, 2.7), donc de présentation finie (I, 6.2), et alors les morphismes $X'(s_i)\to S$ sont de présentation finie (I, 6.2) et il en est par suite de même du morphisme $S'\to S$.

\begin{corollary}[17.16.3]
\label{IV.17.16.3}
\begin{enumerate}
  \item[(i)] Soit $f:X\to S$ un morphisme lisse. Soient $s\in S$, $x$ un point fermé de $X_s$ tel que le corps résiduel $k(x)$ soit séparable sur $k(s)$; alors, dans la conclusion de (17.16.1), on peut prendre $X'$ tel que le morphisme $X'\to U$, restriction de $f$, soit étale.
  \item[(ii)] Soit $f:X\to S$ un morphisme lisse surjectif. Alors, dans les conclusions de (17.16.2), on peut supposer en outre que $g:S'\to S$ soit étale.
\end{enumerate}
\end{corollary}

Il est clair que (ii) se déduit de (i) comme (17.16.2) de (17.16.1), compte tenu de ce que, pour tout $s\in S$, comme $X_s$ est non vide et lisse sur $k(s)$, il existe un point fermé $x\in X_s$ tel que $k(x)$ soit séparable sur $k(s)$ (17.15.10, (iii)). Il suffit donc de prouver (i). On notera que l'anneau $\sh{O}_{X_s,x}$ est alors régulier; si l'on répète la construction faite dans (17.16.1) en prenant pour $(t_i)$ un système régulier de paramètres de $\sh{O}_{X_s,x}$, l'anneau $\sh{O}_{X'_s,x}$ est un corps isomorphe à $k(x)$, donc séparable sur $k(s)$ par hypothèse. La conclusion résulte alors de (17.6.1, $c'$)).

\begin{proposition}[17.16.4]
\label{IV.17.16.4}
Soient $S$ un préschéma quasi-compact et quasi-séparé, $f:X\to S$ un morphisme surjectif localement de présentation finie. Alors il existe une famille finie $(S_i)_{i\in I}$

